\scanpage{39}
géom. du $\infty$-polygône combinatoire épinglé $\Pi_0(\infty)$. Pour ceci, on note que \deleted{la donnée de la} \deleted{dans le cas de $\Pi_0(n)$, et $r$ premier à $n$, on aura la repr. de} (pour $n$ quelconque, fini ou infini) la donnée de la repr.
\[
D_n\mathrel{\makebox[0pt][l]{\makebox[1.35em]{$\rho$}}\unsure{\rule{1.35em}{0.4pt}|\longrightarrow}}\mathrm{Aff}(E)
\]
équivaut aussi à la donnée de
\[
\left\{\begin{aligned}
&\sigma_E,u_E\in\mathrm{Aff}(E)\\
&\text{satisfaisant}\\
&\sigma_E^2= \mathrm{id}_E,\ \sigma_Eu_E\sigma_E^{-1}=u_E
\qquad\text{et si $n$ fini}\\[-2pt]
&\hspace{19em}u_E^n=\mathrm{id}_E
\end{aligned}\right.\tag*{(28)}
\]
tandis que la donnée suppl. $s_0\in E$ et les conditions $(\unsure{2^\circ|3^\circ})$ et $(4^\circ)$ \deleted{\unsure{(ou $3^\circ$, $4^\circ$)}} de la prop 5 \deleted{ne dépendent pas} \deleted{\unsure{survivent}} \inserted{\unsure{sont stables|détaillées} par} le chgt de $u_E$ en $u_E^r$ -- qui a un sens quel que soit $r\in\mathbb Z$, et (pour $n$ fini) continue à vérifier la condition $u_E^n=\mathrm{id}$ (après remplacement de $u_E$ par $u_E^r$). \deleted{\unsure{(Ici il}} Quant à la condition $2^\circ)$ de non align$^{\rm mt}$ de $s_0$, $u^rs_0$ et $u^{2r}s_0$ (ou de tous les $u^{ir}s_0$, $i\in\mathbb Z$, ce qui revient au $=$)
% archive sheet number: 38
